\documentclass{article}
\usepackage{/home/user1/code/tex/sty}
\usepackage{amsfonts}
\usepackage{amsmath}
\usepackage{geometry}
\usepackage{graphicx}
\usepackage{hyperref}
\usepackage{floatrow}
\newcommand{\fig}[3]{
\begin{figure}[H]
\begin{center}
\includegraphics[height=#1]{#2}
\caption{#3}
\label{#2}
\end{center}
\end{figure}}
\usepackage[backend=bibtex]{biblatex}
\addbibresource{bib.bib}
\setlength{\parindent}{0pt}
\geometry{top=1in,bottom=1in,left=1in,right=1in}
\begin{document}
\begin{center}
\textbf{\Large{Problem Set 1}}\\~\\
\begin{tabular}{l l}
Student&Agustin Romero\\
Email&ar1@stanford.edu\\
Instructor&Caterina Vernieri\\
Course&Physics 252\\
Institution&Stanford University\\
Date&\today\\
\end{tabular}
\end{center}
\section*{Problem 1}
The essential experimental tests of chemistry are the following. The Standard Model of Particle Physics should explain
\begin{enumerate}
\item That it supplies a foundation for these experimental facts.
\item The other aspects of the Standard Model are not relevant to chemical interactions.
\item Which of the experimental evidence found in high energy colliders should not be expected to be found in a chemistry lab.
\end{enumerate}
The foundational experimental facts of chemistry are the following.
\begin{enumerate}
\item Oxidation reaction. A chemical loses an electron, and the chemical losing a charge of 1 e.
\e{A \rightarrow A^+ + e^-}
\item Reduction reaction. A chemical gains an electron, and chemical gaining a charge of 1 e.
\e{A + e^- \rightarrow A^-}
\item Redoxreaction. An electron is transferred from one chemical to another.
\e{A + B \rightarrow A^+ + B^-}
\item Van Der Walls forces. Van Der Walls forces can be demonstrated from intermolecular electrostatic forces between molecules.
\item Covalent bonds. Covalent bonds can be demonstrated in the Standard Model from electrodynamics, where electron pairs are shared between atoms.
\item Ionic bonds. Ionic bonds can be demonstrated in the Standard Model from the QED sector, where electrostatic forces attract oppositely charged ions.
\item Alpha radiation. Alpha radiation can be demonstrated in the Standard Model from the QED sector, where the Coulomb force repulses the alpha particle from the nucleus.
\item Beta radiation. Beta radiation can be demonstrated in the Standard Model from the electroweak sector, where W bosons couple to the decaying particle.
\item Gamma radiation. Gamma radiation can be demonstrated in the Standard Model from QED, where high energy particles release some of their kinetic energy as photons.
\item Photon emission. Photon emission in the lab can be modeled as either classical or quantum electromagnetic interactions.
\end{enumerate}
The Standard Model describes the following bosons, only some of which are relevant the 
\begin{enumerate}
\item Photons. Photons are a necessary aspect of chemistry experiments and can be excited at low energies. The QED sector of the Standard Model and its calssical limits provide a good foundatio nfor these particle interactions.
\item Gluons. These are particles whose interaction distance is subatomic, and is almost always not relevant in the laboratory. 
\item W bosons. W bosons are relevant to chemistry in electromagnetic interactions.
\item Z bosons. Z bosons are not to chemistry experiments, becaues their excitation requires high energy, around 90 GeV.
\end{enumerate}
\section*{Problem 2.a.}
The process
\e{mu^+ \rightarrow e^- \nu_e \bar{\nu_e}}
is not possible, because it violates $L_\mu$ number.
\section*{Problem 2.b.}
This is a weak process, where one up quark changes to a down quark and W$^+$ boson.
\fig{4cm}{7.png}{LO diagrams.}
\section*{Problem 2.c.}

\section*{Problem 2.d.}
\section*{Problem 2.e.}
\section*{Problem 2.f.}
\section*{Problem 2.g.}
\section*{Problem 2.h.}
\section*{Problem 2.i.}
\section*{Problem 3.a.}
The lowest order diagrams are
\fig{4cm}{1.png}{LO diagrams.}
\section*{Problem 3.b.}
\fig{4cm}{2.png}{NLO diagrams.}
\section*{Problem 4.}
Label the process
\e{\gamma + p \rightarrow \Delta^+}
as
\e{1 + 2 \rightarrow 3}
Use the invariance of $s$ and find the relativistic energy $\gamma mc^2$ for the proton.
\e{s_i &= s_f
\intertext{Assume the rest frame of the Delta baryon.}
E_3^2 &= (E_1 + E_2)^2 - (\vec{p_1}+\vec{p_2})^2
\intertext{In this frame $\vec{p_1}=-\vec{p_2}$}
E_1 + E_2 &= E_3\\
E_1 + E_2 &= \gamma_3 m_3 c^2 - E_1
\intertext{In this frame the Delta baryon is at rest, $\gamma_3=1$.}
E_2 &= m_3 c^2 - E_1}
Plug in the mass of $\Delta^+$ and use the blackbody energy of a CMB photon.
\e{\boxed{\gamma m_2 c^{-2} = 1232\un{MeV$c^{-2}$} - k_B(2725\un{K})}}
The proton needs essentially 1232 MeV in relativistic energy to excite the $\Delta^+e$ state.
\section*{Problem 5.a.}
Using
\e{m^2=E^3-\vec{p}^2}
\e{p^\mu=(E,p_x,p_y,p_z)}
\e{E=\gamma mc^2}
\e{p=\gamma m \beta}
\e{p^\mu p_\mu = (\sum E_i)^2 - (\sum p_i)^2}
Find the mass $m_\alpha$.
\e{m_\alpha^2 &= (p_1+p_2)^\mu (p_1+p_2)_\mu\\
&=(E_1+E_2)^\mu-(\vec{p_1}+\vec{p_2})^2\\
&=E_1^2+E_2^2+2E_1E_2 - \vec{p_1}^2-\vec{p_2}^2+2\vec{p_2p_3}\\
&=m_1^2+m_2^2+2E_1E_2+2\vec{p_2}\vec{p_3}\\
&=m_1^2+m_2^2+2E_1E_2- 2\gamma_1\gamma_2 m_1 m_2 \beta_1 \beta_2 \cos(\theta)\\
&=m_1^2+m_2^2+2E_1E_2- 2E_1E_2 \beta_1 \beta_2 \cos(\theta)\\
&=\boxed{m_1^2+m_2^2+2E_1E_2(1- v_1 v_2 \cos(\theta))}}
\section*{Problem 5.b.}
Label the process $\alpha \rightarrow 2 + 3$ as before. Assume $m_1$, $m_2$, the momentum $p_1=0.75\un{GeV}$, and $p_2=4.25\un{GeV}$, and the angle $\theta=9^\circ$. Do not assume the mass $m_\alpha$ or the energy $E_1$, $E_2$. Solve for the mass $m_\alpha$, the mass of the $\Delta^+$. Make use of
\e{E=\gamma m}
\e{p=\gamma m \beta}
\e{\lab{91283}\gamma=(1+(\fr{p}{mc})^2)^(1/2)}
\e{\beta=(1-\gamma^{-2})^{1/2}}
Use \eref{91283} to solve for $\gamma$, $\beta$, and $E$ of the $\pi^+$.
\e{\gamma_1 = 5.47}
\e{\beta_1 = 0.98314}
\e{\beta_1 = 763.4\un{MeV}}
Use \eref{91283} to solve for $\gamma$, $\beta$, and $E$ of the proton.
\e{\gamma_2 = 4.64}
\e{\beta_2 = 0.9765}
\e{\beta_2 = 4.352\un{GeV}}
Plug into $m_\alpha^2$ to find the reconstructed mass of the $\Delta^+$.
\e{m_\alpha^2&=m_1^2+m_2^2+2E_1E_2(1- v_1 v_2 \cos(\theta))}
\e{\boxed{m_\alpha=1115.08\un{MeV $c^{-2}$}}}
This result can be compard to the PDG value 1115.683 MeV c$^{-2}$. Next solve for the lifetime using conservation of energy.
\e{E_\alpha &= E_1 + E_2 = 5.1154\un{GeV}}
\e{E=\gamma mc^2}
Using the reconstructed mass $m_\alpha$,
\e{\gamma_\alpha = \fr{E_\alpha}{m_\alpha c^2} = 4.587}
\e{\beta_\alpha = (1-\gamma_\alpha^{-2})^{1/2}= 0.9759}
The average propagation distance $d$ is related to the lifetime $\tau$ by
\e{d=\gamma \beta c \tau}
Solve for $\tau$
\e{\boxed{\tau=2.61\tun{-10}{s}}}
This can be compared to the PDG value 2.632$\tun{-10}{s}$.
\section*{Problem 5.c.}
Label the process $p+p\rightarrow p + p + \bar{p} + \bar{p}$ as $1+2\rightarrow 3+4+5+6$. One initial proton is at rest, and all final state protons are at rest. So let all $\vec{p}=0$ except one of the initial state protons, $\vec{p_2}$.
\e{\vec{p_1}=\vec{p_3}=\vec{p_4}=\vec{p_5}=\vec{p_6}=0}
Then the energy of particle 1 is
\e{E_1=\gamma_1 m_1 c^2 = m_1 c^2 = m_p c^2}
Use the invariance of $s$ to solve for the energy of the energy of the moving proton, particle 2.
\e{(\sum E_i)^2-(\sum \vec{p_i})^2 &= (\sum E_j)^2-(\sum \vec{p_j})^2\\
(E_1+E_2)^2 + \vec{p_2}^2 &= 4m_p^2 c^4 \\
E_1^2 + E_1E_2 + E_2^2 - p_2^2 &= 4m_p^2 c^4\\
E_1^2 + E_1E_2 + m_2^2c^4 &= 4m_p^2 c^4\\
E_1^2 + E_1E_2 &= 5m_p^2 c^4\\
m_p^2c^4 + m_pc^2E_2 &= 5m_p^2 c^4\\
m_pc^2E_2 &= 4m_p^2 c^4}
The energy of the moving proton $E_2$ is 
\e{\boxed{E_2 = 4m_p c^2}}
\section*{Problem 6.a.}
For the first process,
\e{\boxed{m_{\Delta^+}=1232\un{MeV}}}
\e{\boxed{m_{n}=939.59\un{MeV}}}
\e{\boxed{m_{\pi^+}=139.57\un{MeV}}}
\e{\lab{91823n}m_n+m_{pi^+}=1079.13\un{MeV}}
For the second process,
\e{\boxed{m_{\Sigma^0}=1192.642\un{MeV}}}
\e{\boxed{m_{\Lambda^0}=1115.683\un{MeV}}}
\e{\boxed{m_{\gamma}=0\un{eV}}}
For the third process,
\e{\boxed{m_{\pi^+}=139.57\un{MeV}}}
\e{\boxed{m_{\mu^+}=105.65\un{MeV}}}
For the fourth process,
\e{\boxed{m_{\pi^0}=134.976\un{MeV}}}
All processes are kinematically allowed because \fbox{the rest energy of the final state particles is less} than the rest mass of the initial state particles.
\section*{Problem 6.b.}
The Feynman diagram for the first process is
\fig{4cm}{3.png}{Diagram.}
The Feynman diagram for the second process is
\fig{4cm}{4.png}{Diagram.}
The Feynman diagram for the third process is
\fig{4cm}{5.png}{Diagram.}
The Feynman diagram for the fourth process is
\fig{4cm}{6.png}{Diagram.}
\section*{Problem 6.c.}
The first process is a \fbox{weak process}. The second process is an \fbox{electromagnetic process}. The third process is a \fbox{weak process}. The fourth process is an \fbox{electromagnetic process}. Electromagnetic processes typically have a shorter lifetime than weak processes, so process 2 and 4 are shorter than 1 and 3. Process 4 has a larger mass difference than process 2, so process 4 is shorter in time than process 2. Process The process lifetimes are
\begin{enumerate}
\item \fbox{$\tau = 5.57\tun{-24}{s}$}
\item \fbox{$\tau = 7.4\tun{-20}{s}$}
\item \fbox{$\tau = 2.603\tun{-8}{s}$}
\item \fbox{$\tau = 8.52\tun{-17}{s}$}
\end{enumerate}
Which supports the ordering \fbox{4, 2, 3, 1}.
\end{document}
